% Generated draft: requires mathematical and source review.
\documentclass[11pt]{article}
\usepackage[margin=25mm]{geometry}
\usepackage{fontspec}
\setmainfont{Latin Modern Roman}
\usepackage{amsmath,amssymb,mathtools,unicode-math}
\setmathfont{Latin Modern Math}
\usepackage{enumitem,needspace,xurl,hyperref}
\setlength{\parindent}{0pt}
\setlength{\parskip}{5pt}
\setlength{\emergencystretch}{3em}
\newcommand{\N}{\mathbb{N}}
\newcommand{\Z}{\mathbb{Z}}
\newcommand{\Q}{\mathbb{Q}}
\newcommand{\R}{\mathbb{R}}
\newcommand{\C}{\mathbb{C}}
\newcommand{\F}{\mathbb{F}}
\newcommand{\A}{\mathbb{A}}
\newcommand{\PP}{\mathbb{P}}
\newcommand{\E}{\mathbb{E}}
\newcommand{\eps}{\varepsilon}
\newcommand{\dd}{\,\mathrm{d}}
\begin{document}
\section*{UnsolvedMath ID: 30006582}
\textbf{Source problem code:} catalogue entry
\subsection{Definitions and mathematical statement}
\textbf{Formulation unavailable from the supplied passage.} The supplied record is expository and contains no separately formulated mathematical question, conjecture, or theorem to formalize. Its statement status is recorded as \texttt{unrecoverable}. Accordingly, no precise mathematical assertion, domain, hypotheses, boundary conditions, or quantifiers can be recovered without adding information absent from the source record.

The passage describes o-minimality as a model-theoretic framework supporting tame analytic arguments, where ``analytic arguments'' specifically means arguments from real or complex analysis. It mentions P-minimality and C-minimality as approaches to a corresponding framework for non-Archimedean analysis, and states that the contribution works with the more recent notion of \(1\)-h-minimality, described there as more broadly applicable than the preceding notions. The supplied passage does not define these minimality conditions or specify the underlying structures, languages, valued fields, or characteristic assumptions; no formal definitions or comparison theorem are therefore asserted here.

A qualification in the extracted footnote says that everything in the contribution except Example~12 holds for both \(1\)-h-minimality and \(\omega\)-h-minimality. Neither Example~12 nor the assertions covered by this qualification are supplied. The exception does not determine whether Example~12 fails for either variant, or merely requires separate treatment.

The registry labels the entry ``open'', but this is a supplied classification of an expository record, not a verified claim that a particular mathematical proposition is open. The intended mathematical target remains unresolved.
\subsection{Short English statement}
The supplied entry introduces h-minimality as a framework for non-Archimedean analysis but states no precise problem. Its qualification concerning \(1\)-h-minimality, \(\omega\)-h-minimality, and Example~12 cannot be converted into a mathematical assertion without the missing definitions and surrounding text. This statement-only draft requires human mathematical review.
\subsection{Sources}
\begin{itemize}
\item Local problem registry: permanent UnsolvedMath ID \(30006582\), title ``Expository h-Minimality Passage'', source problem code \texttt{catalogue entry}, in the collection ``Oberwolfach Reports Open Problems''. The supplied record gives proposed year 2026 and no separately named proposer.
\item Supplied report citation: \emph{Model Theory and Arithmetic} (2026), Oberwolfach Reports, DOI 10.4171/owr/2026/3, pp.~157--216. Supplied source URL and literature-source link, labelled ``Oberwolfach Report 3/2026'': \url{https://doi.org/10.4171/owr/2026/3}. This report reference is supplied provenance; it has not been independently consulted here.
\item Contribution attribution in the extracted passage: Sebastian Eterovi\'{c}, ``Hensel minimality, p-adic exponentiation and Tate uniformization'', joint work with Floris Vermeulen. The supplied page range covers the report; no contribution-specific page or location for the extracted passage is given.
\item Unresolved reference labels retained from the passage: [16, 17, 20] for applications to the Andr\'{e}--Oort conjecture; [1, 13, 18] for results in functional transcendence; [7] for results regarding the Zilber--Pink conjecture; [10] for P-minimality; [9, 12] for C-minimality; and [5, 6] for \(1\)-h-minimality. Bibliographic entries corresponding to these labels are not supplied.
\item Literature assessment dated \texttt{2026-08-22}: the supplied curation labels the record open, describes the extracted passage as expository rather than a resolved theorem or precise solved problem, and describes the associated non-Archimedean tameness agenda as active. This is supplied dated curation, not an independently verified literature assessment.
\item The supplied statement-verification note says that the official contribution introduces h-minimality and summarizes applications, that the extracted passage contains no separately formulated open problem, and that no clean statement is promoted while the original record is retained publicly for provenance. These are supplied verification notes, not independently established findings here.
\end{itemize}
\end{document}
